\section{Question 3: exécution d'une saisie-contrefaçon}

\subsection{Cas d'un brevet validé en France}

\begin{itemize}
  \item Modalités d’exécution: la saisie-contrefaçon est autorisée par ordonnance sur requête du Tribunal judiciaire de Paris et exécutée en France par un commissaire de justice, conformément à l’article L.615-5 CPI, aux articles R.615-1 à R.615-3 CPI, ainsi qu’à l’article 493 CPC.
  
  \item Champ d’intervention du CPI: le conseil en brevets ne peut pas conduire lui-même la saisie-contrefaçon, laquelle relève exclusivement du commissaire de justice, mais il peut être autorisé à assister aux opérations en qualité de conseil technique si l’ordonnance le prévoit, dans les limites strictes fixées par celle-ci.
  
  \item Missions du CPI: le conseil en brevets peut apporter un éclairage technique lors des constatations, aider à identifier les éléments pertinents au regard des revendications, et assister à la description des produits, procédés ou documents saisis, sans disposer d’un pouvoir autonome d’exécution ou de décision pendant l’opération.
\end{itemize}


\subsection{Cas d'un brevet européen à effet unitaire}

\begin{itemize}
  \item Modalités d’exécution : la conservation des preuves et les descentes sur les lieux sont ordonnées par la JUB sur le fondement de l’article 60 AJUB et mises en œuvre conformément aux règles 192 à 199 RPUPC, notamment par inspection ou prélèvements prévus à la règle 194 RPUPC.
  
  \item Champ d'intervention du CPI: Lors de la descente sur les lieux, le requérant n'est pas présent en personne, mais il 
  peut être représenté par un CPI dont le nom figure dans l'ordonnance de la JUB (article 60(4) AJUB). Toutefois, le CPI ne peut pas exécuter
l’ordonnance de conservation des preuves, celle-ci devant être réalisée par une personne
indépendante désignée par la juridiction et chargée de remettre un rapport écrit (règles 196(4) et
196(5) RPUPC).
\end{itemize}
